\documentclass{article}
\usepackage[UTF8]{ctex}
\usepackage{amsmath,amssymb,amsfonts}
\usepackage{enumitem}
\usepackage{geometry}
\geometry{a4paper, margin=1in}
% === 错题本排版优化（不参与题目标识） ===
\setlength{\parindent}{0pt}
\setlength{\parskip}{0.35em plus 0.10em minus 0.08em}
\setlength{\emergencystretch}{3em}
\setlist{itemsep=0.25em, topsep=0.40em, parsep=0pt}
\allowdisplaybreaks[2]
% === 错题本排版优化结束 ===
\begin{document}
\section*{东北大学2025年数学分析真题}
\section*{73．东北大学2025年数学分析试卷（回忆版）}

\subsection{题目1}
1 ．求 \(\lim _{n \rightarrow \infty} \int_{0}^{1} \ln \left(1+x^{n}\right) \mathrm{d} x\) ．

\subsection{题目2}
2 ．求 \(x^{2}+y^{2} \geq a x(a>0)\) 与 \(z=\sqrt{x^{2}+y^{2}}\) 围成的圆锥体的体积．

\subsection{题目3}
3．求级数 \(\sum_{n=1}^{\infty} n^{3} x^{n}\) 的和函数，其中 \(|x|<1\) 。

\subsection{题目4}
4．求积分 \(I=\iiint_{\Omega} \sqrt{x^{2}+y^{2}+z^{2}} \mathrm{~d} V\) ，其中 \(\Omega: x^{2}+y^{2} \leq z^{2}, x^{2}+y^{2}+z^{2} \leq R^{2}, z \geq 0\) ．

\begin{itemize}

\subsection{题目5}
\end{itemize}

\subsection{题目5}
原题边界“5．”：设 \(f(x)\) 在 \([0,1]\) 上可微，且 \(\left|f^{\prime}(x)\right| \leq M\) ．证明：对任意的正整数 \(n\) ，有
\[
\left|\int_{0}^{1} f(x) \mathrm{d} x-\frac{1}{n} \sum_{i=1}^{n} f\left(\frac{i}{n}\right)\right| \leq \frac{M}{n} .
\]

\subsection{题目6}

\subsection{题目6}
原题边界“6．”：设 \(f(x)\) 在 \([a, b]\) 上可积，且 \(\int_{a}^{b} f(x) \mathrm{d} x=A\) ．证明：
\[
\int_{a}^{b} f(x) e^{f(x)} \mathrm{d} x \int_{a}^{b} \frac{1}{f(x)} \mathrm{d} x \geq(b-a)(b-a+A)
\]

\subsection{题目7}
7．证明函数列 \(f_{n}(x)=\frac{\ln (n x)}{n x^{2}}\) 在 \((1,+\infty)\) 上一致收敛．

\subsection{题目8}
8．设 \(f(x)\) 在 \([a, b]\) 上二阶可导，满足 \(f^{\prime}(a)=f^{\prime}(b)=0\) ．证明：至少存在一点 \(c \in(a, b)\) ，使得

\[
\left|f^{\prime \prime}(c)\right| \geq \frac{4}{(b-a)^{2}}|f(b)-f(a)|
\]

连续．9．证明 \(f(x)=\frac{|\sin x|}{x}\) 在区间 \(J_{1}=(-1,0)\) 和区间 \(J_{2}=(0,1)\) 上均一致连续，但在 \(J_{1} \cup J_{2}\) 上不一致\\

\subsection{题目9}
10．证明反常积分 \(\int_{0}^{+\infty} \frac{1}{1+x^{2} \sin ^{2} x} \mathrm{~d} x\) 发散．
\end{document}
